% 数学ノート集・共通テンプレート v2（積分論2026年度第1部の版面を参考）
% A5判。章内で定義・定理・例・注意等の番号を共通化。
% コンパイル： platex -interaction=nonstopmode math_note_template_v2.tex (2回)
%              dvipdfmx math_note_template_v2.dvi
\documentclass[a5paper,10pt,onecolumn]{jsarticle}
\usepackage[dvipdfmx,top=22mm,bottom=20mm,inner=19mm,outer=18mm,headheight=13pt,headsep=6mm]{geometry}
\usepackage{amsmath,amssymb,amsthm,mathtools}
\usepackage{enumitem,fancyhdr,titlesec}
\usepackage[dvipdfmx,bookmarks=true,colorlinks=false]{hyperref}
\usepackage{pxjahyper}

% ----- 差し替える文書情報 -----
\newcommand{\NoteTitle}{RSAと公開鍵暗号の数学}
\newcommand{\NoteSeries}{数学ノート集}

% ----- 文字・組版 -----
% 参考PDFの質素な明朝＋Computer Modernの印象を引き継ぎながら，
% A5判で窮屈にならないよう本文の行間を少し確保する．
\linespread{1.13}
\setlength{\parindent}{1zw}
\setlength{\parskip}{0pt}
\setlength{\abovedisplayskip}{6pt plus 2pt minus 2pt}
\setlength{\belowdisplayskip}{6pt plus 2pt minus 2pt}
\setlength{\abovedisplayshortskip}{3pt plus 2pt minus 1pt}
\setlength{\belowdisplayshortskip}{4pt plus 2pt minus 1pt}
\setlist[enumerate]{topsep=3pt,itemsep=1pt,parsep=0pt,leftmargin=2.4zw}
\titleformat{\section}{\normalfont\normalsize\bfseries}{\thesection}{0.8zw}{}
\titlespacing*{\section}{0pt}{2.0ex plus .4ex}{.9ex}
\titleformat{\subsection}{\normalfont\normalsize\bfseries}{\thesubsection}{0.7zw}{}
\titlespacing*{\subsection}{0pt}{1.5ex plus .3ex}{.6ex}
\raggedbottom

% ----- 柱とページ番号：ページ上部、罫線なし -----
\fancypagestyle{mathnote}{%
  \fancyhf{}%
  \fancyhead[L]{\footnotesize \NoteSeries\ （\NoteTitle）}%
  \fancyhead[R]{\footnotesize \thepage}%
  \renewcommand{\headrulewidth}{0pt}%
  \renewcommand{\footrulewidth}{0pt}%
}
\pagestyle{mathnote}

% ----- 節内通し番号の定義・定理群 -----
\newtheoremstyle{mathnotetheorem}
 {4.5pt}{4.5pt}{\normalfont}{}%
 {\bfseries}{}{0.65zw}%
 {\thmname{#1}\thmnumber{ #2}\thmnote{（#3）}}
\theoremstyle{mathnotetheorem}
\newtheorem{theorem}{定理}[section]
\newtheorem{definition}[theorem]{定義}
\newtheorem{lemma}[theorem]{補題}
\newtheorem{proposition}[theorem]{命題}
\newtheorem{corollary}[theorem]{系}
\newtheorem{example}[theorem]{例}
\newtheorem{remark}[theorem]{注意}
\renewcommand{\proofname}{証明}
\renewcommand{\qedsymbol}{$\square$}
\makeatletter
\renewenvironment{proof}[1][\proofname]{%
  \par\pushQED{\qed}\normalfont\topsep4\p@\@plus2\p@\relax
  \trivlist\item[\hskip\labelsep\bfseries #1]\ignorespaces
}{\popQED\endtrivlist\@endpefalse}
\makeatother

% ----- 使うノートで適宜拡張する記号 -----
\newcommand{\R}{\mathbb{R}}
\newcommand{\norm}[1]{\lVert #1\rVert}
\newcommand{\abs}[1]{\lvert #1\rvert}

\begin{document}
\section*{到達目標}
有限群・中国剰余定理の知識から，RSAの秘密指数を構成し，
具体的なメッセージを暗号化・復号して一致を検算できるようになる．
代数的な正当性と実用上の暗号安全性は区別する．
前提は合同算術，B\'ezout等式，素数についての基本知識である．


\section{剰余環と単元群}

本稿では合同算術と有限群の基本定理を前提とし，
RSAの復号が正しいことを数学的に示す．
実装上の安全性の問題は，証明可能な代数的事実と区別する．

\begin{definition}
整数 $n\ge2$ に対し，
剰余環 $\mathbb Z/n\mathbb Z$ の\emph{単元群}を
\[
 (\mathbb Z/n\mathbb Z)^\times
 =\{\bar a:\gcd(a,n)=1\}
\]
と定める．
その位数をEuler関数 $\varphi(n)$ と書く．
\end{definition}

\begin{lemma}[単元の判定]\label{lem:units}
$\bar a\in\mathbb Z/n\mathbb Z$ が単元であることと
$\gcd(a,n)=1$ は同値である．
\end{lemma}
\begin{proof}
$\gcd(a,n)=1$ ならBezoutの等式により
$au+nv=1$ となる整数 $u,v$ が存在する．
剰余類へ移せば $\bar a\bar u=\bar1$ なので
$\bar a$ は単元である．
逆に $\bar a\bar u=\bar1$ なら
$au-1$ は $n$ の倍数である．
従って $au+nv=1$ となる整数 $v$ が存在し，
$\gcd(a,n)=1$ である．
\end{proof}

\begin{theorem}[Eulerの定理]\label{thm:euler}
$\gcd(a,n)=1$ なら
\[
 a^{\varphi(n)}\equiv1\pmod n
\]
が成り立つ．
\end{theorem}
\begin{proof}
$\bar a$ は有限群
$(\mathbb Z/n\mathbb Z)^\times$ の元である．
Lagrangeの定理によって $\bar a$ の位数は
$\varphi(n)$ を割るから
$\bar a^{\varphi(n)}=\bar1$ である．
これを合同式で表したものが結論である．
\end{proof}

\begin{corollary}[Fermatの小定理]\label{cor:fermat}
$p$ が素数で $p\nmid a$ なら
$a^{p-1}\equiv1\pmod p$ である．
\end{corollary}
\begin{proof}
$\varphi(p)=p-1$ を定理\ref{thm:euler}に代入する．
\end{proof}

\section{中国剰余定理}

\begin{theorem}[中国剰余定理]\label{thm:crt}
$m,n\ge2$ が互いに素なら
\[
 \mathbb Z/mn\mathbb Z
 \longrightarrow
 \mathbb Z/m\mathbb Z\times\mathbb Z/n\mathbb Z,
 \qquad \bar a\longmapsto(\bar a,\bar a)
\]
は環同型である．
\end{theorem}
\begin{proof}
写像が環準同型であることは剰余類の演算から従う．
核は $m,n$ の両方で割り切れる整数の類である．
$m,n$ が互いに素なので，そのような整数は $mn$ で割り切れ，
核は零である．
従って写像は単射である．
定義域と値域の元の個数はともに $mn$ だから全射でもある．
また構成的には $um+vn=1$ を満たす整数 $u,v$ を取ると，
任意の剰余類 $(\bar b,\bar c)$ の逆像は
\[
 a=bnv+cmu
\]
で与えられる．実際 $a\equiv b\pmod m$，
$a\equiv c\pmod n$ である．
\end{proof}

\begin{corollary}\label{cor:unit-group}
互いに異なる素数 $p,q$ と $n=pq$ に対し
\[
 (\mathbb Z/n\mathbb Z)^\times
 \cong\mathbb F_p^\times\times\mathbb F_q^\times,
 \qquad
 \varphi(n)=(p-1)(q-1)
\]
が成り立つ．
\end{corollary}
\begin{proof}
環同型は単元を単元へ写し，その逆も単元を単元へ写す．
従って定理\ref{thm:crt}から単元群の同型を得る．
$\mathbb F_p^\times,\mathbb F_q^\times$ の位数は
それぞれ $p-1,q-1$ なので位数の等式が従う．
\end{proof}

\section{RSAの鍵生成}

\begin{definition}
互いに異なる素数 $p,q$ を取り $n=pq$ と置く．
\[
 \lambda(n)=\operatorname{lcm}(p-1,q-1)
\]
と定義する．
$\gcd(e,\lambda(n))=1$ を満たす正整数 $e$ を選ぶ．
次に
\[
 ed\equiv1\pmod{\lambda(n)}
\]
となる整数 $d$ を選ぶ．
組 $(n,e)$ を\emph{公開鍵}，
$d$ および素因数 $p,q$ を秘密に保持する量という．
\end{definition}

\begin{lemma}[秘密指数の存在]\label{lem:private}
$\gcd(e,\lambda(n))=1$ なら
$ed\equiv1\pmod{\lambda(n)}$ を満たす整数 $d$ が存在し，
$\lambda(n)$ を法として一意である．
\end{lemma}
\begin{proof}
Bezoutの等式から
$eu+\lambda(n)v=1$ となる整数 $u,v$ が存在する．
$d=u$ と置けば求める合同式を満たす．
二つの解 $d,d'$ があれば
$e(d-d')\equiv0\pmod{\lambda(n)}$ である．
$e$ は $\lambda(n)$ を法として可逆なので
$d-d'\equiv0\pmod{\lambda(n)}$ となり一意である．
\end{proof}

\begin{definition}
平文を剰余類 $\bar m\in\mathbb Z/n\mathbb Z$ とする．
RSA暗号化を
\[
 E_e(\bar m)=\bar m^{\,e}
\]
と定め，復号を
\[
 D_d(\bar c)=\bar c^{\,d}
\]
と定める．
\end{definition}

\section{RSA復号の正当性}

\begin{theorem}[RSAの復号定理]\label{thm:rsa}
互いに異なる素数 $p,q$ を取り $n=pq$ とする．
$\gcd(e,\lambda(n))=1$，
$ed\equiv1\pmod{\lambda(n)}$ とする．
このとき任意の剰余類
$\bar m\in\mathbb Z/n\mathbb Z$ について
\[
 D_d(E_e(\bar m))=\bar m,
 \qquad\text{すなわち}\qquad
 m^{ed}\equiv m\pmod n
\]
が成り立つ．
特に $\gcd(m,n)=1$ は仮定しなくてよい．
\end{theorem}
\begin{proof}
$ed=1+k\lambda(n)$ となる整数 $k\ge0$ が存在する．
$p-1$ は $\lambda(n)$ を割るので
$ed\equiv1\pmod{p-1}$ である．
$m\equiv0\pmod p$ なら
$m^{ed}\equiv m\equiv0\pmod p$ である．
$m\not\equiv0\pmod p$ なら
Fermatの小定理から $m^{p-1}\equiv1\pmod p$ なので
\[
 m^{ed}=m(m^{p-1})^{k\lambda(n)/(p-1)}
 \equiv m\pmod p
\]
となる．
同じ議論で
$m^{ed}\equiv m\pmod q$ を得る．
$p,q$ は互いに素なので中国剰余定理から
$m^{ed}\equiv m\pmod{pq}$ となる．
\end{proof}

\begin{remark}
定理\ref{thm:rsa}は\emph{全ての}剰余類の復号を保証する．
単にEulerの定理を $m$ に適用するだけでは
$\gcd(m,n)=1$ の場合しか直接には扱えない．
$p$ と $q$ をそれぞれ法として場合分けすることが重要である．
\end{remark}

\section{反復二乗法と計算量}

\begin{proposition}[反復二乗法]\label{prop:squaring}
非負整数 $e$ の二進展開
$e=\sum_{j=0}^rb_j2^j$（$b_j\in\{0,1\}$）を用いると，
任意の整数 $a$ に対して
\[
 a^e\bmod n
\]
は高々 $2r+1$ 回程度の剰余乗算で求められる．
\end{proposition}
\begin{proof}
$x_0=a\bmod n$，
$x_{j+1}=x_j^2\bmod n$ と置くと，
帰納法により $x_j\equiv a^{2^j}\pmod n$ となる．
$0\le j<r$ の計算には $r$ 回の二乗が必要である．
次に $b_j=1$ となる $j$ に対してのみ $x_j$ を乗じれば，
\[
 \prod_{j:b_j=1}x_j
 \equiv a^{\sum_jb_j2^j}=a^e\pmod n
\]
となる．
この乗算は高々 $r+1$ 回で済む．
従って剰余乗算の回数は $O(\log(e+1))$ である．
\end{proof}

\begin{proposition}[中国剰余定理による復号]\label{prop:crt-decrypt}
$c\in\mathbb Z$ に対し
$c^d\bmod p$ と $c^d\bmod q$ を求めたとする．
この二つの剰余類から
$c^d\bmod n$ は中国剰余定理によって一意に復元できる．
\end{proposition}
\begin{proof}
$c^d$ の剰余類を $p,q$ それぞれへ写したものは
与えられた二つの値である．
定理\ref{thm:crt}の環同型は全単射なので，
対応する $\mathbb Z/pq\mathbb Z$ の剰余類は一意である．
定理の証明中に与えたBezout係数を用いれば
その剰余類を明示的に計算できる．
\end{proof}

\section{RSAの代数的構造と安全性の区別}

\begin{proposition}[RSA指数写像の置換性]\label{prop:permutation}
定理\ref{thm:rsa}の仮定のもとで
\[
 E_e:\mathbb Z/n\mathbb Z\longrightarrow
 \mathbb Z/n\mathbb Z,\qquad x\longmapsto x^e
\]
は全単射であり，逆写像は $D_d$ である．
\end{proposition}
\begin{proof}
定理\ref{thm:rsa}により
$D_d\circ E_e=\operatorname{id}$ が成立する．
同じ定理は $e$ と $d$ を交換しても成立する：
$de\equiv1\pmod{\lambda(n)}$ であり，
$\gcd(d,\lambda(n))=1$ だからである．
従って
$E_e\circ D_d=\operatorname{id}$ も成り立つ．
よって二つの写像は互いに逆である．
\end{proof}

\begin{example}[小さい具体例]
$p=5$，$q=11$，$n=55$ とする．
$\lambda(n)=\operatorname{lcm}(4,10)=20$ である．
$e=3$ とすれば $d=7$ は
$ed=21\equiv1\pmod{20}$ を満たす．
平文 $m=2$ の暗号文は $c=2^3=8$ であり，
\[
 8^7\equiv2\pmod{55}
\]
となって復号される．
\end{example}
\begin{proof}
$8^2\equiv9$，$8^4\equiv26\pmod{55}$ なので
$8^7=8^4\cdot8^2\cdot8
\equiv26\cdot9\cdot8\equiv2\pmod{55}$
である．
\end{proof}

\begin{remark}
本稿の代数的定理は，公開鍵から秘密情報を計算することが
困難であるとは主張しない．
特にRSAの復号問題の計算量と素因数分解問題の計算量の
厳密な同値性は一般に知られていない．
また，数式上のRSAの正当性は，実際の通信の安全性を
保証するものではない．
決定的な教科書的RSAは選択平文攻撃などへの
安全性を満たさず，実用上は適切な標準化された
符号化・パディング方式が別途必要になる．
\end{remark}

\section{小さな鍵で暗号化と復号を全て計算する}
\begin{example}[鍵生成から復号まで]
$p=5$，$q=11$ を用いる．$n=55$，
$\lambda(n)=\operatorname{lcm}(4,10)=20$ である．
$e=3$ と選べば $\gcd(3,20)=1$ であり，
$3\cdot7=21\equiv1\pmod{20}$ だから秘密指数は
$d=7$ と取れる．公開鍵は $(55,3)$ である．
平文を $m=12$ とすれば暗号文は
\[
 c\equiv12^3=1728\equiv23\pmod{55}.
\]
復号を反復二乗法で行うと
\[
 23^2\equiv34,\qquad
 23^4\equiv34^2\equiv1\pmod{55},
\]
ゆえに
\[
 c^d=23^7\equiv23\cdot34\cdot1
       \equiv12\pmod{55}.
\]
もとの平文が確かに復元された．
\end{example}
\begin{proposition}[単元ではない平文も復号できる]
上の鍵のもとで $m=5$ とする．この平文は
$\mathbb Z/55\mathbb Z$ の単元ではないが，
復号の結果はやはり $5$ になる．
\end{proposition}
\begin{proof}
暗号文は $5^3=125\equiv15\pmod{55}$．
$15^2\equiv5$，$15^4\equiv25$ だから
\[
 15^7\equiv15\cdot5\cdot25\equiv5\pmod{55}.
\]
あるいは中国剰余定理で見ると，法5では
両辺が0，法11では $5^{ed}=5^{21}\equiv5$
がFermatの小定理から成立する．
よって法55でも一致する．
一般の平文に対する復号の正当性は，
平文が単元であることを仮定せず，
素数ごとに場合分けしてから中国剰余定理を使う．
\end{proof}
\section*{到達確認問題と解答}
\noindent\textbf{問題．}
同じ公開鍵 $(55,3)$ で平文 $m=2$ を暗号化し，
$d=7$ で復号せよ．

\medskip\noindent\textbf{解答．}
$2^3\equiv8\pmod{55}$．
$8^2\equiv9$，$8^4\equiv26$，
\[
 8^7\equiv8\cdot9\cdot26\equiv2\pmod{55}.
\]
これも元の平文に一致する．
\section*{数学的な射程}
ここで使った数値は暗号の仕組みを手計算するためのものである．
$n=55$ は容易に素因数分解できるため暗号用途に適さない．
パディングや攻撃モデルを扱わず，この計算から安全性は結論しない．

\end{document}
